% ---
% Agradecimentos
% ---
\begin{agradecimentos}

A conclusão deste trabalho não representa apenas o encerramento de uma etapa acadêmica. Ela carrega o esforço, o cuidado e a confiança de muitas pessoas que me sustentaram até aqui.

Primeiramente, agradeço a Deus, por me conceder força, sabedoria e perseverança para chegar até aqui.

Agradeço à minha mãe, Regenilda de Jesus Lobo, por todo apoio, sacrifício e exemplo de força. Muitas vezes abriu mão de si para cuidar dos filhos, enfrentando o trabalho pesado quando necessário e ensinando, acima de tudo, o valor da integridade. Ao meu pai, José Carlos Batista Oliveira, agradeço pelo apoio constante, pelo esforço para que não me faltasse o necessário, pelas brincadeiras e pela certeza de que posso contar com ele quando preciso.

Às minhas irmãs, agradeço por terem sido figuras maternas em fases diferentes da minha vida. Caiala me acolheu na infância com afeto, cuidado e proteção, ajudando-me a me preparar para o mundo. Carla assumiu esse cuidado na adolescência, incentivando-me a seguir bons caminhos, criando oportunidades e insistindo, do seu jeito, para que eu continuasse em frente. Ao meu irmão Cássio, agradeço por ter sido, mesmo sem notar, um exemplo de generosidade, ensinando-me que a gentileza também aparece nos gestos simples.

Agradeço à Flávia Alessandra, pela parceria e companheirismo nos momentos bons e nos mais desafiadores desta jornada. Sua amizade tornou a vida acadêmica mais leve e é uma das grandes conquistas que levo desse período. Agradeço também a Davi, que, por caminhos inesperados, tornou-se uma referência de incentivo, admiração e valorização da educação, reforçando em mim a ideia de que os estudos podem transformar realidades.

Aos professores da Universidade Estadual de Santa Cruz (UESC) que fizeram parte da minha formação, agradeço pelo conhecimento, pela orientação e pelo incentivo. Em especial, à professora Ximena, pelas contribuições e pelo cuidado no desenvolvimento deste trabalho; ao professor Degas, pela forma leve e descontraída de conduzir suas aulas; e ao meu orientador, professor Jorge Lima, pela orientação desde a residência até esta etapa final, essencial à conclusão deste trabalho e ao meu desenvolvimento profissional.

Por fim, agradeço aos meus amigos Henio, Breno e Thiago, e a todas as pessoas que, mesmo não citadas nominalmente, tiveram importância nessa trajetória.

\end{agradecimentos}
% ---
